% Minimal note: the given feedback controller, and how it is implemented in the augmentation
% training loop. The controller itself is not our design (ruleOfThumb.m, read-only FP model).
% Compile: pdflatex controller-implementation.tex   (twice, for the equation references)
\documentclass[11pt]{article}

\usepackage[a4paper,margin=2.4cm]{geometry}
\usepackage[T1]{fontenc}
\usepackage[utf8]{inputenc}
\usepackage{lmodern}
\usepackage{amsmath,amssymb,mathtools,bm}
\usepackage{microtype}
\usepackage{enumitem}
\usepackage{titlesec}
\usepackage{booktabs}
\usepackage[hidelinks]{hyperref}

\titleformat{\section}{\large\bfseries}{\thesection}{0.7em}{}
\titlespacing*{\section}{0pt}{1.6ex plus .2ex}{0.8ex}

\newcommand{\R}{\mathbb{R}}
\newcommand{\diag}{\operatorname{diag}}
\newcommand{\Cfb}{C_{\mathrm{fb}}}
\newcommand{\Cnorm}{C_{\mathrm{n}}}
\newcommand{\wb}{\omega_{\mathrm{b}}}
\newcommand{\fms}{f_{\mathrm{ms}}}
\newcommand{\ufb}{u_{\mathrm{fb}}}
\newcommand{\udata}{u_{\mathrm{d}}}
\newcommand{\umodel}{u_{\mathrm{m}}}
\newcommand{\ydata}{y_{\mathrm{d}}}
\newcommand{\ymodel}{y_{\mathrm{m}}}
\newcommand{\Yop}{Y_{\mathrm{op}}}
\newcommand{\Tu}{T_{u}}
\newcommand{\Ty}{T_{y}}

\setlist{nosep, leftmargin=1.4em}
\pagestyle{empty}

\begin{document}

\begin{center}
  {\Large\bfseries Implementing the feedback controller in the training loop}\\[0.3em]
  {\normalsize Dirk van den Berg}
\end{center}
\vspace{0.3em}

\noindent The trajectory records were generated with a feedback controller in the loop, so the
same controller has to be present when the model is simulated during training. This note gives
the controller (Section~1), wraps it around the first principles model (Section~2), rearranges
that loop into the form actually used (Section~3), and states how it is placed in the training
code (Section~4).

\medskip\noindent\emph{Notation.}\quad $C$ and $K$ are the damping and stiffness matrices of the
baseline, so the controller and its design gain take the symbols $\Cfb$ and $\kappa_j$.
$C_{\mathrm{d}}$ is the output matrix of the discrete model. $\Tu = \diag(u_{\mathrm{std}})$ and
$\Ty = \diag(y_{\mathrm{std}})$ are the input and output normalisations.

\section{The controller (given)}

The controller is not our design. It is \texttt{ruleOfThumb.m} of the read-only first principles
model, applied per channel to the stage-frame design plant
$\mathrm{sys}(\Yop) = P^{\top} G(M_{\mathrm{op}}(\Yop), C_{\mathrm{damp}}, K)\, P$, frozen at the
record's operating point. With $\wb = 2\pi f_{\mathrm{bw}}$ and $f_{\mathrm{bw}} = 100$~Hz:
\begin{align}
  \Cnorm(s) &=
  \underbrace{\frac{s + \wb/6}{s}}_{\texttt{int}} \cdot
  \underbrace{\frac{s + \wb/3}{s + 3\wb}}_{\texttt{leadlag}} \cdot
  \underbrace{\frac{10\,\wb}{s + 10\,\wb}}_{\texttt{lowpass}} ,
  \label{eq:cnorm}\\[1mm]
  \kappa_j &= \frac{1}{\bigl|\,\mathrm{sys}_{jj}(i\wb)\,\Cnorm(i\wb)\,\bigr|} ,
  \qquad j = 1,2,3 ,
  \label{eq:kappa}\\[1mm]
  \Cfb(z) &= \diag\bigl(\kappa_1\Cnorm,\ \kappa_2\Cnorm,\ \kappa_3\Cnorm\bigr)
  \Bigr|_{\,s \,=\, \frac{2}{T_s}\frac{z-1}{z+1}} .
  \label{eq:cfb}
\end{align}
Only the diagonal entries $\mathrm{sys}_{jj}$ are used, so $\Cfb$ is diagonal in the stage frame,
and the plant enters only through the three numbers $|\mathrm{sys}_{jj}(i\wb)|$, so $\Cfb$ can be
rebuilt at any $\Yop$ at negligible cost. What the implementation needs is its state-space
realisation, three states per channel and biproper:
\begin{equation}
  \Cfb : \quad
  x_{c,k+1} = A_c\,x_{c,k} + B_c\,e_{c,k} ,
  \qquad
  \ufb{}_{,k} = C_c\,x_{c,k} + D_c\,e_{c,k} ,
  \qquad
  x_c \in \R^{9} ,\quad D_c \neq 0 .
  \label{eq:cfbss}
\end{equation}

\section{Wrapping it around the first principles model}

The model is discrete time in normalised coordinates, and its output has no feedthrough:
\begin{equation}
  x_{k+1} = f(x_k, u_k) ,
  \qquad
  y_k = h(x_k) = C_{\mathrm{d}}\,x_k + b ,
  \qquad
  \frac{\partial h}{\partial u} = 0 ,
  \qquad
  x \in \R^{8},\; u,y \in \R^{3} .
  \label{eq:model}
\end{equation}
Closing \eqref{eq:cfbss} around \eqref{eq:model} is the direct implementation: the controller is
driven by the tracking error of the \emph{model} output against the reference, and its force is
injected at the model input alongside the multisine,
\begin{equation}
  \begin{aligned}
    \ymodel{}_{,k} &= h(x_k) , \\
    e_{c,k}        &= r_k - \Ty\,\ymodel{}_{,k} , \\
    \ufb{}_{,k}    &= C_c\,x_{c,k} + D_c\,e_{c,k} , \\
    u_k            &= \Tu^{-1}\bigl(\ufb{}_{,k} + \fms{}_{,k}\bigr) , \\
    x_{k+1}        &= f(x_k,\, u_k) , \\
    x_{c,k+1}      &= A_c\,x_{c,k} + B_c\,e_{c,k} .
  \end{aligned}
  \label{eq:direct}
\end{equation}
The controller was designed in physical units, so $\Ty$ takes the model output to metres and
$\Tu^{-1}$ takes the resulting force back to normalised input. The order of the six lines is not a
choice. By \eqref{eq:cfbss} the controller is biproper, so $\ufb{}_{,k}$ needs $e_{c,k}$ and the
model output must be formed first; by \eqref{eq:model} the model output does not depend on $u_k$,
so the loop closes without an algebraic loop, at one model evaluation per step.

\section{The residual form, which is what is used}

Equation \eqref{eq:direct} needs the reference $r$ and the injected multisine $\fms$ of each
record. Both are stored, so \eqref{eq:direct} is runnable, but both then have to be carried
through the batching alongside the input and output. They can be removed. The record was
generated by the same controller around the machine, so writing the machine loop and the model
loop \eqref{eq:direct} side by side,
\begin{equation}
  \udata = \Cfb\bigl(r - \ydata\bigr) + \fms ,
  \qquad
  \umodel = \Cfb\bigl(r - \ymodel\bigr) + \fms ,
  \label{eq:twoloops}
\end{equation}
and subtracting, with $r$ and $\fms$ identical in both and $\Cfb$ linear, they cancel:
\begin{equation}
  \boxed{\;
    \umodel{}_{,k} \;=\; \udata{}_{,k} \;+\; \Cfb\bigl(\ydata - \ymodel\bigr)_k \;}
  \label{eq:residual}
\end{equation}
The controller now filters the difference between the recorded and the modelled output, and its
force is a correction on the recorded input. Nothing but $\udata$ and $\ydata$ is required, which
is exactly what the loader returns. Realised in the same order as \eqref{eq:direct}:
\begin{equation}
  \begin{aligned}
    \ymodel{}_{,k} &= h(x_k) , \\
    e_{c,k}        &= \Ty\bigl(\ydata{}_{,k} - \ymodel{}_{,k}\bigr) , \\
    \ufb{}_{,k}    &= C_c\,x_{c,k} + D_c\,e_{c,k} , \\
    u_k            &= \udata{}_{,k} + \Tu^{-1}\,\ufb{}_{,k} , \\
    x_{k+1}        &= f(x_k,\, u_k) , \\
    x_{c,k+1}      &= A_c\,x_{c,k} + B_c\,e_{c,k} ,
  \end{aligned}
  \qquad
  \begin{aligned}
    &x_{\tau|\tau} = \mathrm{encoder}(\cdot) , \\[0.5mm]
    &x_{c,\tau|\tau} = 0 .
  \end{aligned}
  \label{eq:recursion}
\end{equation}
The input mean does not appear, because $\Tu^{-1}\ufb$ is an increment on $\udata$ rather than an
absolute input.

\section{The controller is a separate training block}

Equation \eqref{eq:recursion} steps two subsystems in sequence. It is not the joined form, which
would carry $z = [\,x^{\top}\ x_c^{\top}\,]^{\top} \in \R^{17}$ and the single map
\begin{equation}
  z_{k+1} =
  \begin{bmatrix}
    f\Bigl(x_k,\; \udata{}_{,k} + \Tu^{-1}\bigl(C_c x_{c,k}
      + D_c \Ty(\ydata{}_{,k} - h(x_k))\bigr)\Bigr) \\[2mm]
    A_c\,x_{c,k} + B_c\,\Ty\bigl(\ydata{}_{,k} - h(x_k)\bigr)
  \end{bmatrix} .
  \label{eq:joined}
\end{equation}
In the separate form the trainable parameters enter only through the model,
$\theta \in \{f, h\}$, while $(A_c, B_c, C_c, D_c)$ are constant. Two consequences:
\begin{enumerate}
  \item The model state stays at $n_x = 8$ rather than $17$ and, with $x_{c,\tau|\tau} = 0$, the
        encoder never has to reconstruct a controller state. The Jacobian of the first principles
        block is unchanged, since $\Cfb$ acts as a linear filter with fixed coefficients.
  \item $\Cfb$ may differ per record. It is rebuilt at each record's $\Yop$ through
        \eqref{eq:kappa}, nine distinct instances across the 22 records, selected inside the batch
        by a per-sample record index. The joined form \eqref{eq:joined} would have to be rebuilt
        for every record.
\end{enumerate}

\medskip\noindent\textbf{The one approximation.}\quad Everything above is an identity except
$x_{c,\tau|\tau} = 0$. In the residual form the controller filters $\ydata - \ymodel$, which does
not exist before the truncated window opens, because the model trajectory starts at $\tau$. Zero
is the unique value for which the correction in \eqref{eq:residual} vanishes when the model is
exact. What it costs is integral memory: $\Cnorm$ has a pole at $s = 0$, so $\Cfb$ has a pole at
$z = 1$ and never forgets DC, while the validation free run accumulates $x_c$ over the whole
record. That asymmetry is measured, not argued.

\end{document}
